% =====================================================================
% Neuer Abschnitt: Experimentelle Untersuchung der universellen Kodierung
% Einzufügen in gen-db.tex direkt nach dem Abschnitt "Universelle Kodierung
% biologischer Strukturen" (\label{sec:universell}) und vor "Schluss".
%
% Benötigt keine zusätzlichen Pakete (booktabs, array, graphicx, enumitem sind
% in gen-db.tex bereits geladen).
%
% Stand: Das Experiment (src/experiment_search_universal.py) ist vollständig
% beschrieben und vorbereitet; die Messwerte stammen aus einem Lauf, der noch
% aussteht. Alle Messwerte werden daher NICHT im Text behauptet, sondern aus
% den vom Skript erzeugten Dateien eingebunden (\input der Tabellen,
% \includegraphics der Diagramme). Fehlt die Datei, erscheint ein roter Hinweis,
% und das Dokument lässt sich trotzdem setzen.
% =====================================================================

% --- Hilfsmakros ------------------------------------------------------
% Zeitstempel des Messlaufs: nach dem Lauf anpassen (Dateiname in src/results/)
\newcommand{\ucstamp}{ZEITSTEMPEL}
\newcommand{\ucresults}{../src/results}

% Roter Platzhalter für Aussagen, die erst nach dem Messlauf belegt werden können
\newcommand{\ucoffen}[1]{\textcolor{accentred}{\textbf{[#1]}}}

% Tabelle des Skripts einbinden, falls vorhanden
\newcommand{\uctabfile}{search_universal_\ucstamp_tabellen.tex}
\newcommand{\ucinputtables}{%
	\IfFileExists{\ucresults/\uctabfile}%
	{\input{\ucresults/\uctabfile}}%
	{\par\ucoffen{Messlauf ausstehend: \texttt{\expandafter\detokenize\expandafter{\uctabfile}} fehlt}\par}}

% Diagramm einbinden, falls vorhanden: \ucplot{Datei}{Breite}{Beschriftung}{Label}
\newcommand{\ucplot}[4]{%
	\begin{figure}[!ht]
		\centering
		\IfFileExists{\ucresults/#1}%
		{\includegraphics[width=#2\textwidth]{\ucresults/#1}}%
		{\fbox{\parbox{0.8\textwidth}{\centering\ucoffen{Messlauf ausstehend: \texttt{\detokenize{#1}}}}}}
		\caption{#3}
		\label{#4}
	\end{figure}}

% =====================================================================
\section{Experimentelle Untersuchung der universellen Kodierung}
\label{sec:univ_experiment}

Der vorangehende Abschnitt~\ref{sec:universell} ist rein formal. Er beweist, dass sich jede endliche attributierte relationale Struktur verlustfrei als Stapel von Mengenfolgen kodieren lässt (Satz~\ref{satz:univ_verlustfrei}), dass der Vergleich zweier Stapel in erwarteter Linearzeit möglich und bis auf einen konstanten Faktor optimal ist (Sätze~\ref{satz:univ_vergleich} und \ref{satz:univ_untere}) und dass ein invertierter Paarindex das exakte Suchergebnis ohne Einzelvergleiche liefert, dessen Aufwand von der Trefferlast $H$ und nicht von der Datenbankgröße $N$ abhängt (Sätze~\ref{satz:univ_index} und \ref{satz:univ_index_zeit}, Korollar~\ref{kor:univ_unabhaengig}). Offen blieb, ob diese Aussagen in einer realen Implementierung mit Datenbank, Prozess-Pool und Speicherverwaltung tragen. Die Sätze sind Aussagen über Rechenschritte; gemessen wird Wanduhrzeit, und diese enthält Datenbankzugriffe, Serialisierung, Prozesskommunikation und die Kosten des Interpreters. Dieser Abschnitt beschreibt das Experiment, das die Lücke zwischen Beweis und Messung schließt. Er folgt dem Aufbau der früheren Experimente (Abschnitte~\ref{sec:parallel}, \ref{sec:backends} und \ref{sec:multiomics}) und trennt deshalb strikt zwischen dem, was das Experiment prüfen \emph{soll}, dem Versuchsaufbau und den Ergebnissen.

Das Experiment ist das im Abschnitt~\ref{sec:univ_umsetzung} angekündigte Experiment des zweiten Schritts. Es wird durch \texttt{src/experiment\_search\_universal.py} gesteuert und vergleicht bei gleichen Anfragen und gleichen Daten die Indexsuche mit Einzelvergleichen in Python und mit der nativen Bibliothek.

\paragraph{Stand der Auswertung.} Der Versuchsaufbau, die Metriken und die Auswertung sind in diesem Abschnitt vollständig festgelegt. Die Messwerte werden aus den vom Skript erzeugten Dateien (\texttt{search\_universal\_<zeit>.json}, \texttt{\_tabellen.tex} und den Einzeldiagrammen \texttt{plot27\_universal\_*.pdf} bis \texttt{plot45\_universal\_*.pdf}) eingebunden. Wo ein Messwert fehlt, steht im Text ein roter Platzhalter. Der Abschnitt behauptet bewusst keine Zahlen, die nicht aus einem dokumentierten Lauf stammen.

%--------------------------------------------------------------------
\subsection{Fragestellung und Hypothesen}
\label{sec:ucx_hypothesen}

Das Experiment beantwortet vier Fragen, die sich aus den Sätzen des Abschnitts~\ref{sec:universell} ergeben:
\begin{enumerate}[label=(\arabic*)]
	\item \emph{Korrektheit:} Liefern Kern, Paarindex und die bisherigen Multi-Omics-Implementierungen dieselben Entscheidungen, und stimmt das Suchergebnis mit dem von Einzelvergleichen überein?
	\item \emph{Laufzeit:} Wie verhält sich die Ende-zu-Ende-Suchzeit der Indexsuche zu der der Einzelvergleiche in Python und mit der nativen Bibliothek, und aus welchen Phasen setzt sie sich zusammen?
	\item \emph{Skalierung:} Hängt die Suchzeit der Indexsuche, wie Korollar~\ref{kor:univ_unabhaengig} vorhersagt, von der Trefferlast $H$ statt von der Datenbankgröße $N$ ab, während die Einzelvergleiche mit $N$ wachsen?
	\item \emph{Kosten des Index:} Wie groß sind Speicherbedarf und Einfügezeit des Paarindex, und wachsen sie, wie Satz~\ref{satz:univ_index_zeit} sagt, linear mit der Belegung der gespeicherten Stapel?
\end{enumerate}

Aus den Sätzen werden Hypothesen abgeleitet. Jede Hypothese nennt ein Kriterium, an dem sie scheitern kann; eine Hypothese, die keine Messung widerlegen könnte, wäre keine.

\begin{table}[!ht]
	\centering
	\small
	\caption{Hypothesen des Experiments, Herleitung und Widerlegungskriterium.}
	\label{tab:ucx_hypothesen}
	\begin{tabular}{@{}>{\raggedright\arraybackslash}p{0.9cm}>{\raggedright\arraybackslash}p{5.6cm}>{\raggedright\arraybackslash}p{3.6cm}>{\raggedright\arraybackslash}p{4.2cm}@{}}
		\toprule
		Nr. & Aussage & Herleitung & Widerlegt, wenn \\
		\midrule
		H1 & Der Kern \texttt{universal\_core} entscheidet auf Zufalls-, eingebetteten, identischen und umgekehrten Paaren, in beiden Modi und auf Schichtauswahlen wie \texttt{multiomics\_python} (und C++). & Bemerkung~\ref{bem:univ_matrix}: Matrixstapel ist die Multi-Omics-Relation. & mindestens eine abweichende Entscheidung. \\
		H2 & Die Indexsuche liefert dieselbe Treffermenge, Gegenrichtung und Klassifikation wie Einzelvergleiche über alle Stapel, auch nach Löschen und erneutem Einfügen. & Satz~\ref{satz:univ_index}. & mindestens eine abweichende Treffermenge oder Klasse. \\
		H3 & Alle Varianten finden in jeder Anfrage und jedem Modus dieselben Treffer; alle eingebetteten Netzwerke mit bekannter Lösung werden gefunden. & H1, H2 und Konstruktion der Ground Truth. & abweichende Treffermenge oder Trefferquote unter $1$. \\
		H4 & Die Anfragezeit der Indexsuche im Speicher ist eine affine Funktion der Trefferlast $H$ mit hohem Bestimmtheitsmaß und wächst mit $N$ nur so stark wie $H$. & Satz~\ref{satz:univ_index_zeit}. & lineare Anpassung an $H$ schlecht ($R^2$ niedrig) oder Anfragezeit wächst bei konstantem $H$ mit $N$. \\
		H5 & Die Zeit der Einzelvergleiche wächst etwa linear mit $N$ (Exponent $b\approx1$ in $t\sim N^b$), die der Indexsuche deutlich flacher ($b<1$), solange $H$ langsamer als $N$ wächst. & Korollar~\ref{kor:univ_unabhaengig} mit Häufigkeitsannahme. & Exponent der Indexsuche nahe $1$ bei selektiven Anfragen. \\
		H6 & Der Speedup der Indexsuche gegenüber den Einzelvergleichen ist bei selektiven Anfragen (lange Anfragen, mehrere aktive Schichten) groß und sinkt bei unselektiven Anfragen. & Satz~\ref{satz:univ_index_zeit}: Kosten $O(H)$; $H$ wächst mit der Häufigkeit der Paare. & Speedup hängt nicht von der Selektivität ab. \\
		H7 & Im Modus \emph{unabhängig} ist die Indexsuche schneller als im Modus \emph{kohärent}, weil keine Verifikation anfällt. & Satz~\ref{satz:univ_index}, zweite Formel. & Zeitverhältnis kohärent/unabhängig $\le1$ bei vergleichbarer Treffermenge. \\
		H8 & Paarzahl und Wörterbuchgröße des Index wachsen linear mit der Zahl der Netzwerke; die Einfügezeit je Netzwerk ist konstant in $N$. & Satz~\ref{satz:univ_index_zeit}, Speicher- und Pflegekosten. & überlineares Wachstum, steigende Einfügezeit je Netzwerk. \\
		H9 & Die Einzelvergleiche mit nativer Bibliothek sind im Mikrobenchmark schneller als in Python, der Ende-zu-Ende-Gewinn ist aber kleiner, weil Laden und Übergabe dominieren. & Befunde der Abschnitte~\ref{sec:backends} und \ref{sec:multiomics}. & Ende-zu-Ende-Speedup $\ge$ Mikro-Speedup. \\
		\bottomrule
	\end{tabular}
\end{table}

Zwei der Hypothesen formulieren bewusst eine Grenze der Aussage. H5 und H6 gelten nur, solange die Häufigkeitsannahme des Korollars~\ref{kor:univ_unabhaengig} trägt. Bei sehr kurzen Anfragen mit wenigen aktiven Schichten kommt jedes Paar in einem großen Teil der gespeicherten Strukturen vor; dann ist $H=\Theta(N)$, und die Indexsuche liest große Listen, ohne Einzelvergleiche zu sparen. In diesem Fall wäre der Index nicht schneller, und das Experiment soll genau diesen Übergang sichtbar machen (Abschnitt~\ref{sec:ucx_selektivitaet}).

%--------------------------------------------------------------------
\subsection{Untersuchte Varianten}
\label{sec:ucx_varianten}

Verglichen werden vier Varianten mit demselben Schema, denselben Daten und derselben Trefferbedingung. Die Referenz für alle Speedups ist \texttt{python}.

\begin{itemize}
	\item[\texttt{index}] \textbf{Paarindex (SQL).} \texttt{crud.search\_universal} liest die Listen der Tabelle \texttt{universal\_pairs} (Primärschlüssel über Schema, Schicht, Art, Paar und Netzwerk) ohne Prozess-Pool. Im Modus \emph{unabhängig} ist das Ergebnis der Schnitt der Trefferlisten (Satz~\ref{satz:univ_index}), im Modus \emph{kohärent} werden die Kandidaten zusätzlich verifiziert.
	\item[\texttt{python}] \textbf{Einzelvergleiche im Kern.} Alle gespeicherten Stapel des Schemas mit ausreichender Länge werden geladen und im Prozess-Pool mit \texttt{universal\_core.compare\_stacks} einzeln verglichen, ohne Index. Dies entspricht dem Vorauswahlschritt der bisherigen Pipeline (Länge $\ge$ Länge der Anfrage) und ist die Referenz.
	\item[\texttt{mo\_python}] \textbf{Bisherige Multi-Omics-Pipeline in Python.} \texttt{crud.search\_multiomics} mit der Implementierung \texttt{multiomics\_python}: Vorfilter über Schichtnamen und Knotenzahl, Einzelvergleiche im Pool.
	\item[\texttt{cpp}] \textbf{Dieselbe Pipeline mit der nativen Bibliothek.} Die Klasse \texttt{MultiOmics} aus \texttt{csubgraph} wird über \texttt{CSUBGRAPH\_LIB\_PATH} geladen. Vor jeder Messreihe prüft das Skript, dass in den Worker-Prozessen tatsächlich die native Bibliothek aktiv ist; ein stiller Rückfall auf Python würde den Vergleich verfälschen und bricht den Lauf ab.
\end{itemize}

Die Varianten \texttt{python}, \texttt{mo\_python} und \texttt{cpp} nutzen den Prozess-Pool, \texttt{index} nicht. Damit vergleicht das Experiment ausdrücklich Verfahren in ihrem jeweiligen Betriebsmodus und nicht einzelne Rechenkerne: Der Gewinn der Indexsuche enthält auch, dass sie keine Prozesse kommunizieren lässt. Die Mikro-Benchmarks (Abschnitt~\ref{sec:ucx_micro}) trennen den Kern deshalb vom Betriebsmodus.

%--------------------------------------------------------------------
\subsection{Versuchsaufbau}
\label{sec:ucx_aufbau}

\subsubsection{Schema und Kodierung}

Das Experiment verwendet ein Matrixschema (\texttt{universal\_schema.matrix\_schema}, lokale Koordinaten) mit $L_\Sigma$ Schichten (Standard $5$) aus dem Pool Genom, Transkriptom, Proteom, Metabolom, Epigenom und Lipidom. Die Kantenwahrscheinlichkeiten der Zufallsschichten sind $0{,}12$, $0{,}25$, $0{,}30$, $0{,}20$, $0{,}15$ und $0{,}22$. Jedes Netzwerk besitzt alle Schichten des Schemas; dadurch bilden alle Varianten dieselbe Datenbasis desselben Schemas, was Voraussetzung der Sätze des Abschnitts~\ref{sec:universell} ist. Der kodierte Stapel besteht aus einer Existenzschicht, die jede Entität auf sich selbst abbildet, und den In-Nachbarschaften der Matrizen (Bemerkung~\ref{bem:univ_matrix}). Ein Netzwerk wird gleichzeitig als universelle Struktur (Tabellen \texttt{universal\_structures}, \texttt{universal\_components} und \texttt{universal\_pairs}) und als Multi-Omics-Netzwerk (Tabellen \texttt{omics\_*}) abgelegt; beide Darstellungen stammen aus denselben Matrizen. Vor dem Lauf prüft das Skript an Größen von $2$ bis $63$ Knoten, dass die schnelle Erzeugung der Stapel mit der offiziellen Kodierung \texttt{encode\_matrices} und mit \texttt{stack\_from\_matrices} übereinstimmt; bei einer Abweichung bricht es ab.

\subsubsection{Daten und Ground Truth}

Die Datenbank wird um $20.000$ reproduzierbar aus dem Seed ($42$) erzeugte Zufallsnetzwerke des Typs \texttt{universal\_experiment} mit $4$ bis $40$ Knoten ergänzt. Vorhandene andere Netzwerke werden nicht verändert. Zu jeder eingebetteten Anfrage werden zusätzlich Netzwerke mit bekannter Lösung erzeugt:
\begin{itemize}
	\item eine identische Kopie (\texttt{exact}),
	\item je drei \emph{kohärent} eingebettete Netzwerke: zwei benachbarte Spalten der Anfrage werden in allen Schichten an derselben Position und demselben Spaltenpaar in das größere Netzwerk kopiert,
	\item je drei nur \emph{unabhängig} eingebettete Netzwerke (ab zwei Schichten): Position und Spaltenpaar werden je Schicht verschoben, sodass die Anfrage meist nur im Modus unabhängig enthalten ist.
\end{itemize}
Die eingebetteten Netzwerke haben höchstens $12$ Knoten mehr als die Anfrage. Für jedes dieser Netzwerke hält das Skript mit der Python-Referenz fest, in welchem Modus die Anfrage tatsächlich enthalten ist; diese Menge ist die \emph{erwartete Treffermenge}. Sie liefert eine Ground Truth zusätzlich zum Vergleich der Varianten untereinander, denn zwei Varianten, die dieselbe falsche Antwort geben, würden sich sonst nicht unterscheiden.

\subsubsection{Anfragen}

Das Anfragegitter besteht aus sieben Knotenzahlen ($4,6,8,12,16,24,32$), vier Schichtzahlen ($1$ bis $4$) und zwei Arten: einer Zufallsanfrage ohne eingebettete Lösung und einer eingebetteten Anfrage. Es ergeben sich $7\times4\times2=56$ Anfragen. Aktive Schichten sind die Schichten der Anfrage; die übrigen Schichten des Schemas sind leer und nicht aktiv, jede Anfrageschicht enthält mindestens eine Kante. Anfragen haben mindestens vier Knoten. Für Stapel der Länge $1$ ohne Paare unterscheiden sich Indexsuche und \texttt{compare\_stacks} per Definition (Klassifikation des Paarindex); solche Stapel werden nicht verwendet.

\subsubsection{Gleiche Trefferbedingung}

Als Treffer zählen die Ergebnisse, bei denen die Anfrage $Q$ in der gespeicherten Struktur $M_k$ enthalten ist (\texttt{KEEP\_B}, \texttt{EQUAL\_KEEP\_A}, \texttt{EQUAL\_KEEP\_B}, \texttt{IDENTICAL}), also $\pi_I(Q)\sqsubseteq_\mu\pi_I(M_k)$ mit $\mu\in\{\mathrm{ind},\mathrm{coh}\}$. Alle vier Varianten werden an dieser Bedingung gemessen. Treffer anderer Multi-Omics-Netzwerke in der Datenbank, die nicht zum Experiment gehören, zählen bei \texttt{mo\_python} und \texttt{cpp} nicht mit.

\subsubsection{Messprotokoll}

\begin{itemize}
	\item \textbf{Wiederholungen.} Jede Suchmessung wird dreimal wiederholt (Parameter \texttt{--reps}). Je Messreihe wird eine Aufwärmsuche verworfen (\texttt{--warmup}), damit Caches der Datenbank und der Prozess-Pool die Zeiten der ersten Anfrage nicht verfälschen.
	\item \textbf{Reihenfolge.} Die Reihenfolge der Varianten rotiert je Wiederholung, die Reihenfolge der Anfragen wird je Wiederholung aus dem Seed gemischt. Systematische Effekte der Reihenfolge, etwa ein sich erwärmender Prozessor oder gefüllte Datenbank-Caches, verteilen sich so auf alle Varianten.
	\item \textbf{Gemessene Größe.} Gemessen wird die Wanduhrzeit der gesamten Suche (SQL, Vergleiche, Prozesskommunikation, Trefferliste). Zusätzlich werden drei Phasen getrennt: \emph{Laden} (bis zum ersten Vergleich beziehungsweise bis alle Kandidaten geladen sind), \emph{Vergleichen} (Zeit in den Vergleichen) und \emph{Rest} (Umwandlung der Stapel, Aufbau der Treffer). Bei \texttt{index} umfasst das Laden das Lesen der Indexlisten und das Laden der Kandidaten; die Vergleiche klassifizieren dort nur die Kandidaten, im Modus unabhängig sind alle Kandidaten Treffer.
	\item \textbf{Parallelität.} Die Standardausführung verwendet zwei Worker (\texttt{--workers 2}); die Skalierung wird in einem eigenen Teilexperiment untersucht.
	\item \textbf{Störungen.} Der Server darf während des Experiments nicht laufen. Andere Multi-Omics-Netzwerke in der Datenbank, etwa vom Multi-Omics-Experiment, sehen nur \texttt{mo\_python} und \texttt{cpp} (anderes Schema); sie verlangsamen diese beiden Varianten und werden im Bericht als Warnung vermerkt. Vor dem Lauf sollte deshalb \texttt{experiment\_search\_multiomics.py --cleanup-only} ausgeführt werden.
	\item \textbf{Hilfsindex beim Löschen.} Beim Löschen vieler Netzwerke legt das Skript kurzzeitig einen Hilfsindex auf \texttt{universal\_pairs(network\_id)} an, weil die Kaskade von \texttt{biological\_networks} sonst je Netzwerk die ganze Indextabelle durchsucht; der Index wird danach entfernt.
\end{itemize}

Die Versuchsumgebung (Betriebssystem, Python- und NumPy-Version, Prozessor, Kernzahl) schreibt das Skript in die Metadaten der JSON-Datei. Sie ist in Tabelle~\ref{tab:ucx_setup} zu ergänzen, weil die absoluten Zeiten von ihr abhängen und nur die Verhältnisse übertragbar sind.

\begin{table}[!ht]
	\centering
	\small
	\caption{Konfiguration des Experiments (Standardwerte des Skripts, Parameter in Klammern).}
	\label{tab:ucx_setup}
	\begin{tabular}{@{}>{\raggedright\arraybackslash}p{5.2cm}>{\raggedright\arraybackslash}p{8.8cm}@{}}
		\toprule
		Parameter & Wert \\
		\midrule
		Zufallsnetzwerke (\texttt{--networks}) & $20.000$, $4$ bis $40$ Knoten (\texttt{--min-nodes}, \texttt{--max-nodes}) \\
		Schema (\texttt{--schema-layers}) & Matrixschema mit $5$ Schichten aus Genom, Transkriptom, Proteom, Metabolom, Epigenom, Lipidom \\
		Anfrageknotenzahlen (\texttt{--query-sizes}) & $4,6,8,12,16,24,32$ \\
		Anfrageschichtzahlen (\texttt{--query-layers}) & $1,2,3,4$ \\
		Anfragearten und Modi & Zufall/eingebettet; kohärent/unabhängig \\
		Eingebettete Netzwerke (\texttt{--plants}) & $1$ identisch, $3$ kohärent, $3$ unabhängig je eingebetteter Anfrage; bis zu $12$ zusätzliche Knoten \\
		Suchmessungen & $56$ Anfragen $\times\,2$ Modi $\times\,4$ Varianten $\times\,3$ Wiederholungen $=1.344$ (mit allen Varianten) \\
		Korrektheit (\texttt{--correctness-pairs}) & $5.000$ Paare, beide Modi; Indexprüfung auf $150$ Stapeln mit $200$ Anfragen \\
		Mikro-Benchmark & Schichtzahlen $1,2,3,4,6,8$; $n_B\in\{8,16,24,32,48,63\}$; Zufall/eingebettet; je $100$ Paare \\
		Speicherexperiment & $N\in\{250,500,1000,2000,5000\}$; $8$ Anfragen; $5$ Wiederholungen der Indexabfrage \\
		Worker-Skalierung (\texttt{--worker-scaling}) & $1,2,3,4$ Worker \\
		Chunk-Größen (\texttt{--chunksizes}) & $1,8,32,128,256,1024,4096$ \\
		Datenbankgrößen (\texttt{--db-sizes}) & $10\,\%$, $25\,\%$, $50\,\%$, $100\,\%$ von $20.000$, also $2.000$, $5.000$, $10.000$, $20.000$ \\
		Teilmenge für Worker, Chunk, Datenbankgröße & $6$ gleichmäßig gewählte Anfragen, $2$ Wiederholungen \\
		Seed & $42$ \\
		Laufzeitumgebung & \ucoffen{aus den Metadaten der JSON-Datei ergänzen} \\
		\bottomrule
	\end{tabular}
\end{table}

%--------------------------------------------------------------------
\subsection{Teilexperimente und Metriken}
\label{sec:ucx_teilexperimente}

Das Skript gliedert den Lauf in sieben Teilexperimente, die mit \texttt{--only} einzeln gewählt werden können. Die ersten drei laufen ohne Datenbank, die übrigen vier gegen PostgreSQL.

\subsubsection{Korrektheit (ohne Datenbank)}
\label{sec:ucx_korrektheit}

Das Teilexperiment prüft H1 und H2 in zwei Teilen.

\emph{Paarweise Prüfung.} Auf $5.000$ Paaren, gleichmäßig auf fünf Arten verteilt (Zufall, kohärent eingebettet, unabhängig eingebettet, identisch, umgekehrt), werden in beiden Modi \texttt{universal\_core.compare\_stacks}, \texttt{multiomics\_python.compare\_layered} und, falls verfügbar, die native Bibliothek verglichen. Schichtzahl ($1$ bis $4$), Größe der kleineren Struktur ($3$ bis $20$ Knoten) und der größeren ($n_A$ bis $40$) werden zufällig gewählt. Zusätzlich werden Invarianten geprüft: Symmetrie (vertauscht man die Strukturen, vertauschen sich \texttt{KEEP\_A} und \texttt{KEEP\_B}, die übrigen Klassen bleiben), kohärent impliziert unabhängig (Lemma~\ref{lem:univ_eigenschaften}), erwartete Ergebnisse der konstruierten Paare sowie Monotonie unter Projektion: weniger aktive Schichten dürfen keinen Treffer verlieren. Je Paar und Modus wird außerdem der Kern auf einer zufälligen Schichtauswahl gegen die Referenz auf den gewählten Matrizen geprüft.

\emph{Indexprüfung im Speicher.} Ein \texttt{PairIndex} über $150$ zufällige gespeicherte Stapel wird gegen Einzelvergleiche über alle Stapel geprüft. Die $200$ Anfragen sind zur Hälfte Fenster gespeicherter Stapel (mindestens ein Treffer) und zur Hälfte zufällige kurze Stapel. Je Anfrage werden beide Modi und beide Richtungen (\texttt{contains}: Anfrage in gespeicherter Struktur; \texttt{contained}: gespeicherte Struktur in der Anfrage) sowie die Klassifikation (\texttt{classify}) gegen \texttt{compare\_stacks} geprüft, auf allen Schichten und auf zufälligen Schichtauswahlen. Anschließend wird ein Fünftel der Stapel aus dem Index entfernt, die Prüfung auf dem Rest wiederholt und die Stapel werden wieder eingefügt; so wird die Pflege des Index (Satz~\ref{satz:univ_index_zeit}) mitgeprüft.

\emph{Metrik.} Zahl der Prüfungen, Zahl der Abweichungen je Art, Anzahl der Invariantenverletzungen. Erwartet wird $0$ (H1, H2). Die Prüfung ist eine differenzielle Stichprobe mit gezielten Grenzfällen und kein Beweis; sie sichert aber die Umsetzung gegen die bewiesenen Sätze ab.

\subsubsection{Mikro-Benchmark (ohne Datenbank)}
\label{sec:ucx_micro}

Der Mikro-Benchmark misst die Zeit je Einzelvergleich für \texttt{universal\_core}, \texttt{multiomics\_python} und die native Bibliothek in Abhängigkeit von Schichtzahl $L\in\{1,2,3,4,6,8\}$, Knotenzahl $n_B\in\{8,\dots,63\}$ der größeren Struktur (die kleinere hat $\max(3,\lfloor n_B/2\rfloor)$ Knoten), Modus und Paartyp (zufällig oder eingebettet). Alle Implementierungen erhalten dieselben, deterministisch aus Seed und Zelle erzeugten Paare; die Umwandlung in Stapel des Kerns erfolgt vor der Messung. Zusätzlich werden die Aufrufkosten an einem trivialen $1\times1$-Vergleich gemessen. Die Zellen mit $n_B>63$ sind nicht vorgesehen, weil die native Bibliothek eine 64-Bit-Darstellung nutzt (Bemerkung~\ref{bem:univ_matrix}); der Kern selbst hat diese Grenze nicht. Der Mikro-Benchmark liefert den Kernvergleich ohne Datenbank und Pool und stützt H9.

\subsubsection{Indexsuche im Speicher}
\label{sec:ucx_memory}

Dieses Teilexperiment prüft H4, H5 und H8 ohne Datenbankrauschen. Der \texttt{PairIndex} wird mit $N\in\{250,500,1000,2000,5000\}$ Zufallsstapeln und allen eingebetteten Netzwerken aufgebaut; die Aufbauzeit, die Zahl der Listeneinträge und die Größe der Wörterbücher werden gemessen. Für acht gleichmäßig aus dem Anfragegitter gewählte Anfragen und beide Modi wird die Indexabfrage fünfmal ausgeführt (Median) und einmal mit der linearen Suche mit Einzelvergleichen (\texttt{universal\_core.contained}) über alle gespeicherten Stapel verglichen. Je Messung werden die Trefferlast $H$, die Zahl der Kandidaten und die Zahl der Treffer festgehalten; zugleich wird geprüft, dass beide Verfahren dieselbe Treffermenge liefern. Ausgewertet werden
\begin{itemize}
	\item die lineare Anpassung der Anfragezeit an $H$ ($t\approx a+b\,H$, Bestimmtheitsmaß $R^2$), getrennt für alle Messungen und für den Modus unabhängig;
	\item der Exponent $b$ der Anpassung $t\sim N^b$ im doppelt logarithmischen Raum für Index und lineare Suche;
	\item die lineare Anpassung der Aufbauzeit an die Zahl der Stapel.
\end{itemize}

\subsubsection{Ende-zu-Ende-Suche}
\label{sec:ucx_search}

Im Suchexperiment werden alle $56$ Anfragen in beiden Modi mit allen gewählten Varianten gegen die $20.000$ Zufallsnetzwerke und die eingebetteten Netzwerke ausgeführt, dreimal wiederholt. Je Messung werden Gesamtzeit, die drei Phasen, Zahl der Kandidaten, Zahl der Treffer, Zahl der exakten Treffer, die erwarteten und gefundenen eingebetteten Netzwerke sowie ein Prüfwert der Treffermenge (SHA-1 der sortierten Netzwerkkennungen) gespeichert. Kennzahlen:
\begin{itemize}
	\item \textbf{Speedup} einer Variante gegenüber der Referenz \texttt{python}: Quotient der Gesamtzeiten, aggregiert als Quotient der Summen und als \emph{geometrischer Mittelwert} der Speedups je Anfrage mit $95\,\%$-Bootstrap-Intervall.
	\item \textbf{Übereinstimmung der Treffer} aller Varianten je Anfrage und Modus über den Prüfwert (H3).
	\item \textbf{Trefferquote der eingebetteten Netzwerke} (Ground Truth): gefundene geteilt durch erwartete Netzwerke je Anfrage, Modus und Variante.
	\item \textbf{Phasenanteile} und \textbf{Skalierung mit der Kandidatenzahl}.
	\item \textbf{Selektivität}: Zahl der Kandidaten und Treffer nach Schichtzahl, Knotenzahl und Modus.
\end{itemize}

\subsubsection{Worker, Chunk-Größe und Datenbankgröße}
\label{sec:ucx_skalierung}

Die drei letzten Teilexperimente laufen im Modus kohärent auf einer Teilmenge von sechs gleichmäßig gewählten Anfragen mit zwei Wiederholungen. \emph{Workers} variiert die Zahl der Worker-Prozesse ($1$ bis $4$) für die Varianten mit Prozess-Pool und vergleicht das Ergebnis mit den Befunden des Abschnitts~\ref{sec:parallel}. \emph{Chunksize} variiert die Chunk-Größe der parallelen Einzelvergleiche ($1$ bis $4096$). \emph{Dbsize} verändert die Zahl der Zufallsnetzwerke in der Datenbank auf $2.000$, $5.000$, $10.000$ und $20.000$ und misst je Stufe die Suchzeit aller Varianten sowie Speicherbedarf (Zeilen, Paare je Netzwerk, Wörterbuch je Netzwerk) und Einfügezeit des Index. Ausgewertet werden die lineare Anpassung der Gesamtzeit an $N$ und der Exponent $b$ in $t\sim N^b$ je Variante (H5, H8); nach der Reihe stellt das Skript den Ausgangszustand der Datenbank wieder her. Das Experiment ist das datenbankgestützte Gegenstück zum Speicherexperiment: Es enthält zusätzlich SQL, Indexlisten in der Datenbank und Prozesskommunikation.

%--------------------------------------------------------------------
\subsection{Auswertung und Reproduzierbarkeit}
\label{sec:ucx_auswertung}

\paragraph{Statistik.} Mittelwerte über Wiederholungen werden mit $95\,\%$-Konfidenzintervallen der $t$-Verteilung angegeben, Quotienten über Anfragen als geometrische Mittelwerte mit Bootstrap-Intervall ($2.000$ Ziehungen), weil Speedups multiplikative Größen sind. Die Anpassungen an $H$ und $N$ werden mit der Methode der kleinsten Quadrate bestimmt und mit $R^2$ angegeben. Die Exponenten $b$ stammen aus der linearen Regression von $\log t$ auf $\log N$ und gelten nur im gemessenen Bereich; $N$ umfasst bei der Datenbankgröße nur vier Stufen, bei dem Speicherexperiment fünf, sodass die Exponenten Schätzungen und keine Konstanten sind.

\paragraph{Ausgabedateien.} Der Lauf schreibt in \texttt{src/results/} mit Zeitstempel: die JSON-Datei mit Metadaten, Anfragen, allen Einzelmessungen und Auswertung (wird während des Laufs fortlaufend aktualisiert, ein Abbruch verliert höchstens einige Sekunden), ein PDF mit den Diagrammen, einen Textbericht, die Tabellen für dieses Dokument (\texttt{\_tabellen.tex}) sowie die Einzeldiagramme \texttt{plot27\_universal\_*.pdf} bis \texttt{plot45\_universal\_*.pdf}. Aus einer vorhandenen JSON-Datei lassen sich PDF, Diagramme, Tabellen und Bericht ohne neuen Messlauf erzeugen (\texttt{--plot-from}).

\paragraph{Aufruf.} Der Aufruf erfolgt im Projektordner, damit die \texttt{.env} gefunden wird:
\begin{lstlisting}[language=bash,caption={Aufruf des Experiments}]
python src/experiment_search_universal.py --quick        # kleiner Vorabtest
python src/experiment_search_universal.py                # vollstaendig, dauert Stunden
python src/experiment_search_universal.py --only correctness micro memory
python src/experiment_search_universal.py --backends index python --networks 5000
python src/experiment_search_universal.py --plot-from src/results/search_universal_<zeit>.json
\end{lstlisting}
Die Experiment-Netzwerke bleiben standardmäßig in der Datenbank, damit der Lauf wiederholbar ist und schneller neu startet; \texttt{--cleanup} entfernt sie am Ende, \texttt{--rebuild-data} erzeugt sie neu. Nach dem Lauf ist in diesem Dokument der Zeitstempel (Makro \texttt{\textbackslash ucstamp}) einzutragen.

%--------------------------------------------------------------------
\subsection{Ergebnisse}
\label{sec:ucx_ergebnisse}

Die folgenden Tabellen und Diagramme werden aus dem Messlauf eingebunden. Die Lesehilfen unter den Abschnitten legen vor der Messung fest, welcher Befund welche Hypothese stützt oder widerlegt; sie enthalten keine Messwerte.

\ucinputtables

\subsubsection{Korrektheit}

Tabelle~\ref{tab:uc_correctness} fasst die paarweise Prüfung je Paartyp zusammen. Die Hypothesen H1 und H2 gelten, wenn alle Abweichungs- und Verletzungszähler $0$ sind. Jede Abweichung ist ein Fehler der Umsetzung oder des Beweises und wird vor allen Laufzeitaussagen analysiert; die Beispiele stehen in der JSON-Datei. Ergebnis: \ucoffen{nach Lauf: Zahl der Paare, Abweichungen, Verletzungen einsetzen}.

\ucplot{plot38_universal_korrektheit.pdf}{0.75}{Korrektheit: Entscheidungen des Kerns gegen \texttt{multiomics\_python} und die native Bibliothek je Paartyp.}{fig:ucx_korrektheit}

\subsubsection{Ende-zu-Ende-Suche}

Tabelle~\ref{tab:uc_summary} stellt je Variante und Modus Gesamtzeit, Phasen, Kandidaten, Treffer und Speedup gegenüber der Referenz zusammen, Tabelle~\ref{tab:uc_recall} die Trefferquote der eingebetteten Netzwerke. H3 gilt, wenn die Trefferprüfwerte aller Varianten je Anfrage und Modus übereinstimmen und die Trefferquote der Ground Truth $1$ beträgt. Ergebnis: \ucoffen{nach Lauf: Speedup \texttt{index} gegen \texttt{python}, geometrischer Mittelwert mit Intervall, getrennt nach Modus}.

\ucplot{plot27_universal_gesamtzeit_speedup.pdf}{0.85}{Gesamtzeit und Speedup der Varianten gegenüber \texttt{python}.}{fig:ucx_gesamtzeit}
\ucplot{plot28_universal_verteilung_antwortzeiten.pdf}{0.85}{Verteilung der Antwortzeiten je Variante.}{fig:ucx_verteilung}
\ucplot{plot29_universal_phasen.pdf}{0.85}{Phasen der Suche: Laden, Vergleichen, Rest.}{fig:ucx_phasen}

Die Phasendiagramme (Abbildung~\ref{fig:ucx_phasen}) beantworten die Frage, wo die verbleibende Zeit der Indexsuche liegt. Nach Abschnitt~\ref{sec:multiomics_search} entfiel bei C++ der größte Teil der Zeit auf das Laden; die Indexsuche lädt im Modus unabhängig nur die tatsächlichen Treffer. H9 erwartet, dass sich die Verhältnisse entsprechend verschieben. Ergebnis: \ucoffen{nach Lauf: Phasenanteile je Variante}.

\subsubsection{Selektivität, Modi, Schichten und Knoten}
\label{sec:ucx_selektivitaet}

\ucplot{plot30_universal_skalierung_kandidaten.pdf}{0.85}{Suchzeit in Abhängigkeit von der Kandidatenzahl.}{fig:ucx_kandidaten}
\ucplot{plot31_universal_schichten.pdf}{0.85}{Suchzeit und Speedup nach Schichtzahl der Anfrage.}{fig:ucx_schichten}
\ucplot{plot32_universal_knoten.pdf}{0.85}{Suchzeit und Speedup nach Knotenzahl der Anfrage.}{fig:ucx_knoten}
\ucplot{plot33_universal_selektivitaet.pdf}{0.85}{Selektivität: Kandidaten und Treffer nach Anfrageart.}{fig:ucx_selektivitaet}
\ucplot{plot34_universal_modi.pdf}{0.85}{Vergleich der Modi kohärent und unabhängig.}{fig:ucx_modi}
\ucplot{plot45_universal_treffer_erwartet.pdf}{0.85}{Gefundene und erwartete Treffer der eingebetteten Netzwerke.}{fig:ucx_treffer}

Diese Abbildungen prüfen H6 und H7. H6 gilt, wenn der Speedup der Indexsuche mit wachsender Knoten- und Schichtzahl der Anfrage (kleineres $H$) zunimmt und für kurze Anfragen mit einer Schicht am kleinsten ist. Ergibt sich für kurze Anfragen ein Speedup unter $1$, so ist dies kein Fehlschlag des Experiments, sondern der in Abschnitt~\ref{sec:ucx_hypothesen} vorhergesagte Fall $H=\Theta(N)$, und die Anwendung des Index ist dort nicht sinnvoll. H7 gilt, wenn die Zeit im Modus unabhängig bei vergleichbarer Treffermenge kleiner ist als im Modus kohärent. Ergebnis: \ucoffen{nach Lauf: Anfragen mit Speedup $<1$ benennen, Übergang beschreiben}.

\subsubsection{Mikro-Benchmark}

\ucplot{plot35_universal_mikro_zeit.pdf}{0.85}{Zeit je Einzelvergleich nach Schichtzahl und Knotenzahl.}{fig:ucx_mikro_zeit}
\ucplot{plot36_universal_mikro_speedup_zufall.pdf}{0.85}{Mikro-Speedup auf zufälligen Paaren.}{fig:ucx_mikro_zufall}
\ucplot{plot37_universal_mikro_speedup_eingebettet.pdf}{0.85}{Mikro-Speedup auf eingebetteten Paaren.}{fig:ucx_mikro_eingebettet}

Tabelle~\ref{tab:uc_micro} und \ref{tab:uc_overhead} liefern die Zeit je Vergleich und die Aufrufkosten. Der Mikro-Speedup der nativen Bibliothek gegenüber dem Kern wird mit dem Ende-zu-Ende-Speedup verglichen; H9 gilt, wenn Letzterer kleiner ist. Ergebnis: \ucoffen{nach Lauf: Mikro-Speedup (geometrischer Mittelwert über Zellen) und Aufrufkosten}.

\subsubsection{Skalierung mit der Datenbankgröße}

\ucplot{plot39_universal_memory_zeit.pdf}{0.85}{Indexsuche gegen lineare Suche im Speicher in Abhängigkeit von $N$.}{fig:ucx_memory_zeit}
\ucplot{plot40_universal_memory_trefferlast.pdf}{0.85}{Anfragezeit der Indexsuche in Abhängigkeit von der Trefferlast $H$ mit linearer Anpassung.}{fig:ucx_memory_h}
\ucplot{plot41_universal_dbsize.pdf}{0.85}{Suchzeit aller Varianten in Abhängigkeit von der Datenbankgröße.}{fig:ucx_dbsize}

Tabelle~\ref{tab:uc_memory} gibt für das Speicherexperiment je $N$ Aufbau- und Anfragezeit, Trefferlast und Speedup an, Tabelle~\ref{tab:uc_dbfit} die Anpassungen in der Datenbank. Abbildung~\ref{fig:ucx_memory_h} prüft H4: Gilt die Anfragezeit der Indexsuche als affine Funktion von $H$ mit hohem $R^2$, so wird Satz~\ref{satz:univ_index_zeit} gestützt. Die Exponenten prüfen H5; ein Exponent der Einzelvergleiche nahe $1$ bestätigt den linearen Aufwand der Referenz (Satz~\ref{satz:univ_vergleich} je Vergleich, $N$ Vergleiche), ein Exponent der Indexsuche deutlich unter $1$ stützt Korollar~\ref{kor:univ_unabhaengig}. Ergebnis: \ucoffen{nach Lauf: Exponenten mit $R^2$, Steigung und Achsenabschnitt der Anpassung an $H$}.

\subsubsection{Speicherbedarf und Einfügezeit}

\ucplot{plot42_universal_speicher.pdf}{0.85}{Speicherbedarf des Paarindex und Einfügezeit je Netzwerk.}{fig:ucx_speicher}

Tabelle~\ref{tab:uc_storage} fasst Zeilenzahl, Paare je Netzwerk, Wörterbuchgröße und Einfügezeit zusammen. H8 gilt, wenn die Paarzahl linear mit der Zahl der Netzwerke wächst ($R^2$ nahe $1$) und die Einfügezeit je Netzwerk nicht mit $N$ wächst. Der Index ist der Preis der schnelleren Suche: Die Zahl der Einträge je Netzwerk ist von der Größenordnung der Länge des Stapels mal der Schichtzahl (Satz~\ref{satz:univ_index_zeit}), und dieser Mehrbedarf ist gegen den Gewinn abzuwägen. Ergebnis: \ucoffen{nach Lauf: Paare je Netzwerk, Wörterbuch je Netzwerk, Einfügezeit je Netzwerk}.

\subsubsection{Worker und Chunk-Größe}

\ucplot{plot43_universal_workers.pdf}{0.85}{Skalierung mit der Worker-Anzahl.}{fig:ucx_workers}
\ucplot{plot44_universal_chunksize.pdf}{0.85}{Einfluss der Chunk-Größe.}{fig:ucx_chunksize}

Die Tabellen~\ref{tab:uc_workers} und \ref{tab:uc_chunksize} ordnen die Befunde in die der Abschnitte~\ref{sec:parallel} und \ref{sec:multiomics} ein. Die Indexsuche läuft ohne Pool und ist von beiden Parametern unabhängig; sie dient als Referenzlinie. Ergebnis: \ucoffen{nach Lauf: Speedup S(p) und Effizienz je Variante, günstigste Chunk-Größe}.

%--------------------------------------------------------------------
\subsection{Diskussion und Grenzen}
\label{sec:ucx_diskussion}

\paragraph{Was jedes Ergebnis bedeuten würde.} Die Auswertung trennt drei Fälle. Bestätigt sich H1 bis H4, so tragen die Beweise des Abschnitts~\ref{sec:universell} bis in die Implementierung: Kern und Index sind korrekt, und die Laufzeitaussage über die Trefferlast ist belegt. Bestätigen sich H1 bis H3, aber nicht H5 und H6, so ist die Häufigkeitsannahme des Korollars~\ref{kor:univ_unabhaengig} für die gewählte Datenverteilung verletzt; die Sätze bleiben richtig, ihre Aussage ist aber für dieses Schema und diese Anfragen praktisch ohne Gewinn. Scheitern H1 bis H3, so liegt ein Fehler vor, den es vor jeder Laufzeitaussage zu beheben gilt. Das Experiment ist so angelegt, dass es alle drei Fälle unterscheiden kann.

\paragraph{Zusammenhang mit der bisherigen Pipeline.} Die Bemerkung zum Verhältnis zur bisherigen Pipeline (Abschnitt~\ref{sec:univ_index}) sagte voraus, dass die Indexsuche im Modus unabhängig nur die tatsächlichen Treffer lädt, während die Pipeline alle Kandidaten lädt. Die Phasenanalyse prüft dies unmittelbar. Zu beachten ist, dass \texttt{mo\_python} und \texttt{cpp} zusätzlich einen Vorfilter über Schichtnamen und Knotenzahl nutzen, den die Variante \texttt{python} in der Länge des Stapels nachbildet; die Kandidatenzahlen der Varianten sind daher vergleichbar, aber nicht gleich.

\paragraph{Grenzen der Untersuchung.}
\begin{itemize}
	\item \textbf{Synthetische Daten.} Die Zufallsnetzwerke haben je Schicht eine feste Kantenwahrscheinlichkeit. Reale Netzwerke sind nicht gleichverteilt: Hubs und häufige Motive lassen $H$ stark anwachsen. Die Aussagen über die Selektivität sind daher nur für diese Verteilung belegt. Die Ground Truth stammt aus konstruierten Einbettungen und ersetzt keine biologische Validierung (Bemerkung~\ref{bem:univ_semantik}).
	\item \textbf{Ein Schema.} Das Experiment prüft ein Matrixschema. Die Kernunabhängigkeit (Korollar~\ref{kor:univ_kern}) wird damit nur für diesen Spezialfall gemessen; Schemata für Hypergraphen, Pfade oder Ontologien sind nicht Teil des Experiments. Die Messungen bestätigen die Kodierung nur indirekt über den Spezialfall des Matrixstapels (Bemerkung~\ref{bem:univ_matrix}).
	\item \textbf{Größenbereich.} Anfragen und Netzwerke haben höchstens $40$ Knoten (Mikro-Benchmark bis $63$), weil die native Bibliothek an die 64-Bit-Darstellung gebunden ist. Das Verhalten bei sehr großen Netzwerken, für die der Kern keine Grenze hat, ist nicht gemessen.
	\item \textbf{Datenbankgröße.} Mit $20.000$ Zufallsnetzwerken liegt die Datenbank weit unter den $10^6$ Netzwerken der Evaluation (Abschnitt~\ref{sec:evaluation}). Die Exponenten gelten im gemessenen Bereich von vier beziehungsweise fünf Stufen.
	\item \textbf{Hardware und Betrieb.} Absolute Zeiten hängen von Prozessor, Kernzahl und Datenbankkonfiguration ab; übertragbar sind nur die Verhältnisse. Der Index wird in PostgreSQL gelesen und hängt vom Cache der Datenbank ab; Messungen mit kalter und warmer Datenbank sind nicht getrennt.
	\item \textbf{Kein Vergleich mit der klassischen Isomorphie.} Das Experiment misst die definierte Relation $\sqsubseteq$ und nicht die klassische Subgraph-Isomorphie (Bemerkung~\ref{bem:univ_np}); ein Vergleich mit VF2 gehört nicht dazu.
	\item \textbf{Kopplung.} Die Messungen der Varianten \texttt{mo\_python} und \texttt{cpp} sind durch fremde Multi-Omics-Netzwerke in der Datenbank beeinflussbar; das Skript vermerkt dies als Warnung.
\end{itemize}

%--------------------------------------------------------------------
\subsection{Fazit}
\label{sec:ucx_fazit}

Das Experiment stellt die formalen Aussagen über die universelle Kodierung auf die Probe, mit vier Varianten bei gleichem Schema, gleichen Daten und gleicher Trefferbedingung, mit einer konstruierten Ground Truth und mit Hypothesen, die an benannten Kriterien scheitern können. Es trennt die Korrektheit (Kern und Index gegen Einzelvergleiche), die Kosten des Einzelvergleichs (Mikro-Benchmark), die Skalierung der Indexsuche mit der Trefferlast und der Datenbankgröße (Speicher- und Datenbankexperiment), die Ende-zu-Ende-Suche mit ihren Phasen sowie Speicherbedarf und Pflegekosten des Index. Welche Aussagen die Messung bestätigt, ist erst nach dem Lauf entscheidbar; die Auswertung ist so festgelegt, dass das Ergebnis nicht nachträglich gedeutet werden kann. Der Abschnitt~\ref{sec:schluss} fasst den Stand zusammen und leitet daraus den Ausblick ab.
